\subsection{Damping and near-quadratic tracking}
\label{sec:at-tracking}

We now compare the actual autonomous closure to the dense flow.  All
constants in this subsection are independent of width, order and physical
time.  Fix an initialization in \eqref{eq:at-good-event}.  Denote the full
dense backward carriers by
$k^{(L)}_{D,a}=w_D$ and
$k^{(\ell)}_{D,a}=(W_D^{(\ell+1)})^\top\delta^{(\ell+1)}_{D,a}$
for $\ell<L$.  The preceding Gaussian argument gives, simultaneously at
integer cutoffs $M\ge1$,
\begin{align}
 H_n(M,t)&=\sum_{\ell=1}^L\max_a
 \frac{\|k^{(\ell)}_{D,a}(t)
                \mathbf1_{|k^{(\ell)}_{D,a}(t)|>M}\|_2}{\sqrt n},\nonumber\\
 Z_n(M)&=\int_0^\infty\rho_D(t)H_n(M,t)\dd t
       \le Ce^{-cM^2}+a_n,\qquad a_n\xrightarrow{\Pr}0.
 \label{eq:at-tracking-tail}
\end{align}
Here the indicator acts coordinatewise, and $a_n$ depends only on the dense
reference.  The physical bounds imply $H_n\le C$.  Put
\[
 D=\sup_{t\ge0}d_n(\widehat\theta(t),\theta_D(t)),\qquad
 \epsilon=\int_0^\infty e_E(t)\dd t,\qquad
 Q_r=\int_0^\infty\|\widehat r(t)-r_D(t)\|_m\dd t.
\]
All three quantities are finite before making any comparison, by the
all-order fitting and speed bounds.  They depend on $n,q$; this dependence
is suppressed only in this proof.

\paragraph{One-reference damping.}
Forward subtraction through the bounded operators and slopes gives
\[
 \max_{\ell,a}\frac{\|\widehat z_a^{(\ell)}-z_{D,a}^{(\ell)}\|_2
                 +\|\widehat h_a^{(\ell)}-h_{D,a}^{(\ell)}\|_2}{\sqrt n}
 \le C d_n(t).
\]
Indeed, at each hidden link subtract as
$(\widehat W-W_D)h_D+\widehat W(\widehat h-h_D)$ and use the feature
RMS and operator bounds.  In backward subtraction the only problematic
product is a changed gate multiplied by a dense carrier.  Splitting at $M$
gives
\begin{equation}
 \frac{\|[g_\ell(\widehat z)-g_\ell(z_D)]\od k_D\|_2}{\sqrt n}
 \le jM\frac{\|\widehat z-z_D\|_2}{\sqrt n}
       +2s\frac{\|k_D\mathbf1_{|k_D|>M}\|_2}{\sqrt n}.
 \label{eq:at-cutoff-product}
\end{equation}
The other terms are a bounded gate times
$\widehat W^\top(\widehat\delta-\delta_D)$ and a bounded gate times
$(\widehat W-W_D)^\top\delta_D$.  At the top the first term is instead
the readout difference.  Descending through the fixed depth therefore gives
\begin{equation}
 \max_{\ell,a}\frac{\|\widehat\delta_a^{(\ell)}
                       -\delta_{D,a}^{(\ell)}\|_2}{\sqrt n}
 \le C[(1+M)d_n(t)+H_n(M,t)].
 \label{eq:at-backward-difference}
\end{equation}
The factor $M$ is added at each layer; it is not multiplied by another
cutoff when propagating the next backward difference.

Expanding the products in the tangent Gram
\eqref{eq:at-tangent-gram}, Cauchy--Schwarz yields the same bound for
$\|\widehat\Gamma-\Gamma_D\|_{\rm op}$.  The sample factor $1/m$
makes a matrix whose entries are bounded by $B/m$ have operator norm at
most $B$.  For $u=\widehat r-r_D$ the exact residual difference equation is
\[
 \dot u=-2\widehat\Gamma u
        -2(\widehat\Gamma-\Gamma_D)r_D+\widehat J E.
\]
The preserved Gram gap and $\|\widehat J E\|_m\le CY e_E$ imply
\[
 D^+\|u\|_m\le-\lambda\|u\|_m
       +C\rho_D[(1+M)d_n+H_n]+Ce_E.
\]
At zeros of $u$ this inequality follows by differentiating
$\sqrt{\|u\|_m^2+\zeta^2}$ and then letting $\zeta\downarrow0$.
Integrating and dropping the nonnegative terminal norm gives
\begin{equation}
 \int_0^t\|u\|_m\dd s
 \le C\left[(1+M)\int_0^t\rho_Dd_n\dd s
                  +Z_n(M;t)+\epsilon(t)\right],
 \label{eq:at-integrated-residual}
\end{equation}
where a final argument $t$ means that an integral stops at $t$.

Subtract the readout updates as
$u\widehat h+r_D(\widehat h-h_D)$.  In a hidden update use the three
differences of the residual, backward response and forward response.
The rank-one Frobenius identity and \eqref{eq:at-backward-difference} give
\[
 d_n(t)\le C\epsilon(t)+C\int_0^t\|u\|_m\dd s
                 +C\int_0^t\rho_D[(1+M)d_n+H_n]\dd s.
\]
Insert \eqref{eq:at-integrated-residual}.  On each terminal interval,
replace its nondecreasing source by its terminal value and apply the
integrating factor for $C(1+M)\rho_D$.  Since $\int_0^\infty\rho_D\le CY$,
then taking the supremum gives
\begin{equation}
 D\le A_M(\epsilon+Z_n(M)),\qquad A_M=Ce^{KM}\ge1,
 \qquad Q_r\le C[(1+M)D+Z_n(M)+\epsilon].
 \label{eq:at-damping}
\end{equation}
Only the dense reference supplied the carrier tail in this argument.

\paragraph{A smoother history used only in the proof.}
Let $\operatorname{clip}_M$ act coordinatewise by truncation to $[-M,M]$.
At the actual closure state define auxiliary backward fields
\begin{align}
 \delta_{a}^{(L),M}
   &=g_L(\widehat z_a^{(L)})\od\operatorname{clip}_M(\widehat w),\nonumber\\
 \delta_{a}^{(\ell),M}
   &=g_\ell(\widehat z_a^{(\ell)})\od
      \operatorname{clip}_M\bigl((\widehat W^{(\ell+1)})^\top
                                  \delta_a^{(\ell+1),M}\bigr).
 \label{eq:at-proof-clipping}
\end{align}
These fields are not evolved or used by the closure.  Clipping contracts
Euclidean distance and decreases magnitudes, so their RMS norms are at
most $CY$.  At the top, the almost-everywhere chain rule gives
\[
 \frac{\|\dot\delta_a^{(L),M}\|_2}{\sqrt n}
 \le jM\frac{\|\dot{\widehat z}_a^{(L)}\|_2}{\sqrt n}
       +s\frac{\|\dot{\widehat w}\|_2}{\sqrt n}.
\]
At a lower layer the corresponding upper bound is
\[
 jM\frac{\|\dot{\widehat z}_a^{(\ell)}\|_2}{\sqrt n}
 +s\|\dot{\widehat W}^{(\ell+1)}\|_{\rm op}
                \frac{\|\delta_a^{(\ell+1),M}\|_2}{\sqrt n}
 +s\|\widehat W^{(\ell+1)}\|_{\rm op}
                \frac{\|\dot\delta_a^{(\ell+1),M}\|_2}{\sqrt n}.
\]
Lipschitz functions composed with absolutely continuous coordinates are
absolutely continuous, so this calculation also applies when $g_\ell$ or
clipping is not classically differentiable at every point.  Downward
induction and \eqref{eq:at-speeds} give
$\|\dot\delta_a^{(\ell),M}\|_2/\sqrt n\le C(1+M)\widehat\rho$.
Clipping the readout is essential here: unbounded activations give an RMS
bound on $w$, not a width-independent coordinate bound.

Set $c=\widehat r/\widehat\rho$ and $b_a^{(\ell),M}=c_a\delta_a^{(\ell),M}$,
extended by zero over the prefix.  The zero readout makes the join
continuous.  Since $\|\dot c\|_m\le C$ and $m$ is fixed,
\begin{equation}
 \frac{\|\dot b_a^{(\ell),M}(t)\|_2}{\sqrt n}
 \le C[1+(1+M)\widehat\rho(t)].
 \label{eq:at-auxiliary-speed}
\end{equation}
Freeze this proof history after physical time $T$; denote the frozen history
by $b_T^M$.  Its difference from $b^M$ has normalized clock $L^2$ norm
at most $Ce^{-\kappa T/2}$, because its amplitude is bounded and its
support has length at most $\int_T^\infty\widehat\rho$.
For any current endpoint $A=\tau(t)$, the weighted derivative energy obeys
\begin{align}
 \frac1n\int_0^A\xi(A-\xi)\|(b_T^M)'(\xi)\|_2^2\dd\xi
 &\le\frac{\sup_t\tau(t)}{\kappa n}
           \int_0^{\min(t,T)}\|\dot b^M(s)\|_2^2\dd s\nonumber\\
 &\le C[T+(1+M)^2].
 \label{eq:at-weighted-auxiliary}
\end{align}
Indeed $A-\tau(s)\le\int_s^\infty\widehat\rho
\le\widehat\rho(s)/\kappa$ cancels the inverse clock speed.
The remaining integral is bounded using \eqref{eq:at-auxiliary-speed} and
$\int_0^\infty\widehat\rho^2<\infty$.  On each finite interval before
freezing the residual is positive, so the frozen history is in $H^1$.
The weighted Legendre estimate \eqref{eq:at-weighted-tail} and projection
contraction now give
\begin{equation}
 \left[\frac1n\int_0^A\|(I-\Pi_q^A)b_a^{(\ell),M}\|_2^2\dd\xi\right]^{1/2}
 \le C\left[\frac{1+M+\sqrt T}{q}+e^{-\kappa T/2}\right].
 \label{eq:at-backward-tail}
\end{equation}

\paragraph{Returning to the actual backward history.}
Compare \eqref{eq:at-proof-clipping} with the dense backward recursion.
At the top, insert the dense clipped readout: clipping contraction controls
the readout difference, the changed gate costs $CMd_n$, and the omitted
dense tail is a term of $H_n$.  At every lower layer the clipped-input
difference is bounded by the next auxiliary/dense backward difference plus
$CY\|\widehat W-W_D\|_F$.  The changed gate and clipping tail have the
same bounds as at the top.  Induction, and then
\eqref{eq:at-backward-difference}, yield
\[
 \max_{\ell,a}\frac{\|\widehat\delta_a^{(\ell)}
                             -\delta_a^{(\ell),M}\|_2}{\sqrt n}
 \le C[(1+M)d_n+H_n].
\]
Multiplying by the same $c_a$ and integrating in the closure's own clock gives
\begin{equation}
 \left[\frac1n\int_0^A\|b_a^{(\ell)}-b_a^{(\ell),M}\|_2^2\dd\xi\right]^{1/2}
 \le C(1+M)D+C\sqrt{Z_n(M)+Q_r}.
 \label{eq:at-auxiliary-mismatch}
\end{equation}
For the tail term, use $H_n\le C$ and
$|\widehat\rho-\rho_D|\le\|\widehat r-r_D\|_m$ to obtain
$\int\widehat\rho H_n^2\le C Z_n(M)+C Q_r$.
There is no assumed bound on the ratio of the two residuals.

The forward history tail is at most $C/q$ by
\eqref{eq:at-forward-energy} and \eqref{eq:at-weighted-tail}.
Use this factor, \eqref{eq:at-backward-tail},
\eqref{eq:at-auxiliary-mismatch}, and the exact absolute-defect inequality
\eqref{eq:at-defect-product}.  Increasing its terminal time to infinity
gives
\begin{equation}
 \epsilon\le C\left[
 \frac{1+M+\sqrt T}{q^2}+\frac{e^{-\kappa T/2}}q
 +\frac{(1+M)D}{q}+\frac{\sqrt{Z_n(M)+Q_r}}q\right].
 \label{eq:at-quadratic-source}
\end{equation}
The bounds are uniform in the terminal time and the accumulated absolute
defect increases to $\epsilon$, which justifies this limit.

\paragraph{Absorbing the remaining feedback.}
Put $U=\epsilon+Z_n(M)$.  By \eqref{eq:at-damping},
$D\le A_MU$ and $Z_n(M)+Q_r\le C(1+M)A_MU$.  Consequently
\[
 U\le Z_n(M)+C\left[\frac{1+M+\sqrt T}{q^2}
                       +\frac{e^{-\kappa T/2}}q\right]
       +\frac{C(1+M)A_M}{q}U
       +\frac{C\sqrt{(1+M)A_M}}q\sqrt U.
\]
If $C(1+M)A_M/q\le1/4$, absorb the linear term and apply
$h\sqrt U\le U/4+h^2$ to the last term.  Multiplication by $A_M$ gives
\begin{equation}
 D\le CA_MZ_n(M)
  +CA_M\left[\frac{1+M+\sqrt T}{q^2}
                              +\frac{e^{-\kappa T/2}}q\right]
  +\frac{C(1+M)A_M^2}{q^2}.
 \label{eq:at-cutoff-conclusion}
\end{equation}
Set $u=q^{-2}+a_n$.  For small $u$, choose an integer
$M\le C+C\sqrt{\log(e+1/u)}$ making $Ce^{-cM^2}\le u$, and
$T\le C\log(e+1/u)$ making $e^{-\kappa T/2}\le\sqrt u$.
Then $Z_n(M)\le Cu$.  Since $u\ge q^{-2}$, the cutoff is at most
$C+C\sqrt{\log(e+q)}$, so the absorption condition holds above one
order threshold independent of width and of $a_n$.  Each term of
\eqref{eq:at-cutoff-conclusion} is bounded by
\begin{equation}
 C\Phi(q^{-2}+a_n),\qquad
 \Phi(u)=u\exp\{K\sqrt{\log(e+1/u)}\},\qquad\Phi(0)=0.
 \label{eq:at-phi-bound}
\end{equation}
Enlarging $K$ absorbs the logarithmic factors.  The common physical bound
on $D$ handles the finitely many smaller orders and nonsmall $u$.
Finally $\Phi(u)/u$ is decreasing, hence
$\Phi(v+w)\le\Phi(v)+\Phi(w)$: multiply its value at $v+w$ separately
by $v$ and $w$.  Set $b_n=C\Phi(a_n)$ and enlarge $K$ once more to
bound $\Phi(q^{-2})$ by $\omega(q)$.  This proves
\eqref{eq:alltime-parameters}.  Continuity of $\Phi$ at zero gives
$b_n\to0$ in probability, uniformly in the choice of order.

\subsection{Whole-input prediction and population convergence}

The full first-layer Frobenius bound in \eqref{eq:at-good-event}, together
with \eqref{eq:at-speeds}, controls $\|W^{(1)}\|_F/\sqrt n$ for both
paths at every time.  Linear activation growth and forward induction give,
for any $x\in\R^d$,
\begin{align}
 \frac{\|h^{(\ell)}(t,x)\|_2}{\sqrt n}
       &\le C(1+\|x\|_2/\sqrt d),\nonumber\\
 |\widehat f(t,x)-f_D(t,x)|
       &\le C(1+\|x\|_2/\sqrt d)d_n(t).
 \label{eq:at-whole-input}
\end{align}
For the second inequality, subtract the forward recursions as before and
expand the readout pairing into readout difference and feature difference.
The constant term is needed for activations with nonzero offsets.
Taking the time supremum before integrating proves the closure-to-dense
version of \eqref{eq:alltime-predictions}; taking a supremum on a bounded
input set proves its absolute-error counterpart.

To identify the dense target on all inputs, add the first-layer Gaussian
row coordinates and any finite collection of input probes to the dense
reference programs.  These are passive observations and never enter the
training residuals.  The finite-program limit and finite-time reference
comparison already proved give convergence in probability of their
predictions on every compact time interval.  The global slope and weight
bounds imply, for both dense and closure predictors,
\begin{align}
 |f(t,x)-f(t,x')|&\le C\|x-x'\|_2/\sqrt d,\nonumber\\
 |\partial_t f(t,x)|&\le Ce^{-\kappa t}(1+\|x\|_2/\sqrt d).
 \label{eq:at-predictor-equicontinuity}
\end{align}
For the second inequality, differentiate the forward recursion using
\eqref{eq:at-speeds}, then differentiate the final readout pairing.
The bounds hold uniformly in $n,q$ on the common events.

A finite net of any bounded input set and the first inequality extend
dense probe convergence to the entire set.  The integrated second
inequality bounds the remaining variation after time $T$ by
$Ce^{-\kappa T}(1+\|x\|_2/\sqrt d)$, uniformly in width.  Choose $T$
large, then take the width limit on $[0,T]$.  This proves all-time dense
prediction convergence uniformly on bounded input sets to the uniquely
constructed dense population predictor $f_\infty$.
For a general fixed test law with finite second moment, truncate to a
bounded input set.  The complement contributes at most a constant times
$\int_{\|x\|>R}(1+\|x\|_2/\sqrt d)^2\dd\mu$, which tends to zero.
Thus $d_{n,\mu}:=\mathcal E_\mu(f_{n,D},f_\infty)\to0$ in probability.
The triangle inequality and \eqref{eq:at-whole-input} give
\eqref{eq:alltime-predictions} with
$b_{n,\mu}=C_\mu b_n+d_{n,\mu}$, independent of $q$.
This completes the proof of \Cref{thm:alltime}.

For precision about fixed-order limits, compactify time by
$u=e^{-\kappa t}\in[0,1]$.  Equations
\eqref{eq:at-whole-input}--\eqref{eq:at-predictor-equicontinuity}
give uniform boundedness and equicontinuity on every compact input set,
including the endpoint $u=0$.  Finite nets and a diagonal selection over
integer input radii prove compactness in local uniform convergence.
The same second-moment majorant upgrades this to the test norm
\eqref{eq:test-metric}.  On events of probability tending to one this
gives tightness of each fixed-order predictor family.  Every subsequential
limit law inherits $\mathcal E_\mu(F_q,f_\infty)\le C_\mu\omega(q)$
almost surely: for any positive margin the finite-width violation
probability tends to zero.  This identifies accuracy of predictor limit
points, without asserting uniqueness of their law or of a population
moment-state ODE.
